% Esquisse d'un programme (cote 119, copy 1): archive cote 119, site page 42 (PDF page 43), typescript page 40.
% typescript. Draft compared: EsquisseFr.pdf (text layer):1449-1509 ; EsquisseFr.pdf (text layer):1521-1539.
% AI-transcribed by Claude Opus 5.5 (claude-opus-5-5, adaptive thinking, effort high), Message Batches API, batch pilot (msgbatch_0189Nhf1xB615YWnztubtLfc), 2026-11-12; not yet checked by eye.
\scanpage{119/42}
\mspage{40}

des habitudes acquises, mais un travail tenace, méticuleux, sans
doute de longue haleine, et porteur au jour le jour de moissons
éloquentes.

Je voudrais encore dire quelques mots sur une réflexion
plus ancienne (fin des années 60 ?), très proche de celle dont il
vient d'être question, inspirée par les idées de Nash, qui m'avaient
beaucoup frappé. Au lieu ici de définir axiomatiquement une notion
de ``théorie modérée'' via la donnée d'une notion de ``partie modérée
de $\mathbb{R}^n$ '' satisfaisant à certaines conditions (de stabilité surtout),
c'est à une axiomatisation de la notion de ``variété lisse'' et du
formalisme différentiable sur de telles variétés que j'en avais,
via la donnée, pour chaque entier naturel $n$, d'un sous-anneau
$\mathbb{A}_n$ de l'anneau des germes de fonctions réelles à l'origine dans $\mathbb{R}^n$.
Ce sont les fonctions qui seront admises pour exprimer les ``change-%
ments de carte'' pour la notion de $\mathbb{A}$-variété correspondante, et il
s'est agit de dégager tout d'abord un système d'axiomes sur le
système $\mathbb{A} = (\mathbb{A}_n)_{n\in\underline{N}}$ qui assure à cette notion de variété une sou-%
plesse comparable à celle de variété $C^\infty$, ou analytique réelle
(ou de Nash). Suivant le type de constructions familières qu'on
tient à pouvoir effectuer dans le contexte des $\mathbb{A}$-variétés, le
système d'axiomes pertinent est plus ou moins réduit, ou riche.
Très peu suffit s'il s'agit seulement de développer le formalisme
différentiel, avec la construction de fibrés de jets, les complexes
de De Rham etc. Si on veut un énoncé du type ``quasi-fini implique
fini'' (pour une application au voisinage d'un point), qui est ap-%
paru comme un énoncé-clef dans la théorie locale des espaces ana-%
lytiques, il faut un axiome de stabilité de nature plus délicate,
dans le ``Vorbereitungssatz'' de Weierstrass (*). Dans d'autres

\rule{4cm}{0.4pt}

(*) Il peut paraître plus simple de dire que les anneaux (locaux)
$\mathbb{A}_n$ sont \uline{henséliens}, ce qui est équivalent. Mais il n'est nullement
clair a priori sous cette dernière forme que la condition en ques-%
tion est dans la nature d'une condition de stabilité, circonstance
importante comme il apparaîtra dans les réflexions qui suivent.

\margin{42}
